\documentclass[12pt,a4paper]{article}

% --- GÓI TIẾNG VIỆT VÀ FONT TIMES NEW ROMAN ---
\usepackage[utf8]{vietnam}
\usepackage{mathptmx}

% --- KHỔ GIẤY VÀ CĂN LỀ ---
\usepackage[a4paper, left=2.5cm, right=2cm, top=2.5cm, bottom=2.5cm]{geometry}

% --- TOÁN HỌC ---
\usepackage{amsmath,amssymb,amsfonts,amsthm}

% --- DANH SÁCH, BẢNG ---
\usepackage{enumitem}
\usepackage{booktabs}
\usepackage{array}

% --- LIÊN KẾT ---
\usepackage{hyperref}
\hypersetup{colorlinks=true, linkcolor=blue, urlcolor=blue}

% --- MÔI TRƯỜNG ĐỀ BÀI ---
\theoremstyle{definition}
\newtheorem{cau}{Câu}

\newcommand{\giai}{\noindent\textbf{\textit{Giải.}}\ }
\newcommand{\dapso}[1]{\noindent\fbox{\textbf{Đáp số:} #1}}

\begin{document}

\begin{center}
    {\large \textbf{PHẦN I. LÝ THUYẾT XÁC SUẤT}}\\[0.3em]
    {\LARGE \textbf{BÀI TẬP CHƯƠNG 1. XÁC SUẤT CỦA BIẾN CỐ}}\\[0.8em]
    {\normalsize Họ và tên: \dotfill\quad MSSV: \dotfill}
\end{center}

\vspace{0.5em}
\hrule
\vspace{1em}

\noindent\textbf{Các công thức sử dụng (theo bài giảng Chương 1):}
\begin{itemize}[leftmargin=1.5em]
    \item Xác suất cổ điển: $P(A) = \dfrac{k}{n} = \dfrac{\text{Số trường hợp thuận lợi cho } A}{\text{Số trường hợp có thể xảy ra}}$.
    \item Công thức cộng: $P(A \cup B) = P(A) + P(B) - P(AB)$; nếu $A, B$ xung khắc thì $P(A \cup B) = P(A) + P(B)$; $P(A) = 1 - P(\overline{A})$.
    \item Xác suất có điều kiện: $P(A|B) = \dfrac{P(AB)}{P(B)}$, \ $P(B) > 0$.
    \item Công thức nhân: $P(AB) = P(A)\,P(B|A)$; nếu $A, B$ độc lập thì $P(AB) = P(A)\,P(B)$.
    \item Công thức xác suất đầy đủ và Bayes với hệ đầy đủ $\{A_i\}_{i=1}^n$:
    \[
    P(B) = \sum_{i=1}^n P(A_i)\,P(B|A_i), \qquad P(A_k|B) = \frac{P(A_k)\,P(B|A_k)}{P(B)}.
    \]
    \item Công thức Bernoulli: $P_n(k, p) = C_n^k\, p^k (1-p)^{n-k}$, \ $k = 0, 1, \dots, n$.
\end{itemize}

\vspace{0.5em}
\hrule
\vspace{1em}

%==================================================================
\begin{cau}
Một cửa hàng có 30 máy tính, trong đó có 20 máy tính do công ty A sản xuất và 10 máy tính do công ty B sản xuất. Một khách hàng đến cửa hàng mua 3 máy tính. Giả sử khả năng được mua của mỗi máy là như nhau. Tính xác suất để khách hàng này mua được 2 máy của công ty A và 1 máy của công ty B.
\end{cau}

\giai
Gọi $H$: ``Khách hàng mua được 2 máy của công ty A và 1 máy của công ty B''.

Phép thử là chọn ngẫu nhiên 3 máy từ 30 máy, các cách chọn là đồng khả năng nên ta dùng định nghĩa xác suất dạng cổ điển.
\begin{itemize}
    \item Số trường hợp có thể xảy ra:
    \[
    n = C_{30}^3 = \frac{30 \cdot 29 \cdot 28}{3!} = 4060.
    \]
    \item Số trường hợp thuận lợi cho $H$: chọn 2 máy trong 20 máy của công ty A \textit{và} chọn 1 máy trong 10 máy của công ty B (quy tắc nhân):
    \[
    k = C_{20}^2 \cdot C_{10}^1 = 190 \cdot 10 = 1900.
    \]
\end{itemize}
Vậy
\[
P(H) = \frac{k}{n} = \frac{C_{20}^2 \cdot C_{10}^1}{C_{30}^3} = \frac{1900}{4060} = \frac{95}{203} \approx 0{,}4680.
\]
\dapso{$P(H) = \dfrac{95}{203} \approx 0{,}4680 = 46{,}80\%$.}

\vspace{1em}
%==================================================================
\begin{cau}
Một hãng nước ngọt thăm dò thị hiếu người tiêu dùng về 2 loại nước giải khát A và B. Trong số 200 người được hỏi có 120 người thích loại A, 100 người thích loại B, 60 người thích cả hai loại A và B. Chọn ngẫu nhiên 1 người trong số được thăm dò. Tính xác suất người này:
\begin{enumerate}[label=\arabic*.]
    \item Thích ít nhất một loại nước giải khát trên.
    \item Không thích loại nước nào cả.
\end{enumerate}
\end{cau}

\giai
Gọi
\begin{itemize}
    \item $A$: ``Người được chọn thích loại A'' $\implies P(A) = \dfrac{120}{200} = 0{,}6$.
    \item $B$: ``Người được chọn thích loại B'' $\implies P(B) = \dfrac{100}{200} = 0{,}5$.
    \item $AB$: ``Người được chọn thích cả hai loại'' $\implies P(AB) = \dfrac{60}{200} = 0{,}3$.
\end{itemize}
\begin{enumerate}[label=\arabic*.]
    \item Biến cố ``Thích ít nhất một loại'' là $A \cup B$. Vì $A$ và $B$ không xung khắc ($AB \ne \emptyset$), áp dụng công thức cộng cho hai biến cố tùy ý:
    \[
    P(A \cup B) = P(A) + P(B) - P(AB) = 0{,}6 + 0{,}5 - 0{,}3 = 0{,}8.
    \]
    \item Gọi $K$: ``Người được chọn không thích loại nào cả'' $\implies K = \overline{A} \cap \overline{B}$.\\
    Theo công thức De Morgan: $K = \overline{A} \cap \overline{B} = \overline{A \cup B}$. Do đó
    \[
    P(K) = 1 - P(A \cup B) = 1 - 0{,}8 = 0{,}2.
    \]
\end{enumerate}
\dapso{1. $0{,}8$;\quad 2. $0{,}2$.}

\vspace{1em}
%==================================================================
\begin{cau}
Một hộp phấn có 10 viên trắng, 6 viên đỏ. Chọn ngẫu nhiên không hoàn lại 2 viên (mỗi lần lấy 1 viên).
\begin{enumerate}[label=\arabic*.]
    \item Tính xác suất viên thứ nhất màu trắng.
    \item Tính xác suất viên thứ hai màu đỏ, biết rằng viên thứ nhất màu trắng.
\end{enumerate}
\end{cau}

\giai
Gọi $T_i$: ``Viên phấn lấy lần thứ $i$ màu trắng''; $D_i$: ``Viên phấn lấy lần thứ $i$ màu đỏ'' ($i = 1, 2$). Hộp có tổng cộng $10 + 6 = 16$ viên.
\begin{enumerate}[label=\arabic*.]
    \item Ở lần lấy thứ nhất, 16 viên đồng khả năng được chọn, trong đó có 10 viên trắng:
    \[
    P(T_1) = \frac{10}{16} = \frac{5}{8} = 0{,}625.
    \]
    \item Cần tính xác suất có điều kiện $P(D_2|T_1)$.

    \textbf{Cách 1 (thu hẹp không gian mẫu):} Khi $T_1$ đã xảy ra (đã lấy ra 1 viên trắng, không hoàn lại), hộp còn 15 viên gồm 9 trắng và 6 đỏ. Do đó
    \[
    P(D_2|T_1) = \frac{6}{15} = \frac{2}{5} = 0{,}4.
    \]
    \textbf{Cách 2 (theo định nghĩa):} Số cách lấy có thứ tự 2 viên là $16 \cdot 15 = 240$; số cách lấy viên 1 trắng, viên 2 đỏ là $10 \cdot 6 = 60$. Suy ra $P(T_1 D_2) = \dfrac{60}{240} = \dfrac{1}{4}$, nên
    \[
    P(D_2|T_1) = \frac{P(T_1 D_2)}{P(T_1)} = \frac{1/4}{5/8} = \frac{2}{5} = 0{,}4.
    \]
\end{enumerate}
\dapso{1. $\dfrac{5}{8} = 0{,}625$;\quad 2. $\dfrac{2}{5} = 0{,}4$.}

\vspace{1em}
%==================================================================
\begin{cau}
Một lớp có 20 sinh viên, trong đó có 17 sinh viên giỏi toán; 8 sinh viên giỏi anh; trong đó có 5 sinh viên giỏi cả hai môn. Với giả thiết chọn được một sinh viên giỏi toán, tính xác suất để sinh viên đó cũng giỏi anh.
\end{cau}

\giai
Chọn ngẫu nhiên 1 sinh viên trong lớp. Gọi
\begin{itemize}
    \item $T$: ``Sinh viên được chọn giỏi toán'' $\implies P(T) = \dfrac{17}{20}$.
    \item $A$: ``Sinh viên được chọn giỏi anh'' $\implies P(A) = \dfrac{8}{20}$.
    \item $AT$: ``Sinh viên được chọn giỏi cả hai môn'' $\implies P(AT) = \dfrac{5}{20}$.
\end{itemize}
Cần tính $P(A|T)$.
\begin{itemize}
    \item \textbf{Cách 1 (theo định nghĩa):}
    \[
    P(A|T) = \frac{P(AT)}{P(T)} = \frac{5/20}{17/20} = \frac{5}{17} \approx 0{,}2941.
    \]
    \item \textbf{Cách 2 (thu hẹp không gian mẫu):}
    \[
    P(A|T) = \frac{\text{Số sinh viên giỏi toán và giỏi anh}}{\text{Số sinh viên giỏi toán}} = \frac{5}{17} \approx 0{,}2941.
    \]
\end{itemize}
\dapso{$P(A|T) = \dfrac{5}{17} \approx 0{,}2941$.}

\vspace{1em}
%==================================================================
\begin{cau}
Tung 2 đồng xu cân đối và đồng chất. Gọi $A$: biến cố đồng xu thứ nhất xuất hiện mặt ngửa; $B$: biến cố đồng xu thứ hai xuất hiện mặt ngửa; $C$: biến cố 2 đồng xu xuất hiện mặt giống nhau. Xác định không gian mẫu và biểu diễn $A, B, C$ dưới dạng tập hợp. Từ đó tính $P(A), P(B), P(C)$ và chứng minh $A, B, C$ đôi một độc lập.
\end{cau}

\giai
Ký hiệu $N$: mặt ngửa, $S$: mặt sấp; kết quả $NS$ nghĩa là đồng xu thứ nhất ngửa, đồng xu thứ hai sấp.
\begin{itemize}
    \item \textbf{Không gian mẫu:}
    \[
    \Omega = \{NN,\ NS,\ SN,\ SS\}, \qquad n = 4.
    \]
    Do hai đồng xu cân đối, đồng chất nên 4 biến cố sơ cấp là đồng khả năng, mỗi biến cố sơ cấp có xác suất $\dfrac{1}{4}$.
    \item \textbf{Biểu diễn các biến cố và tính xác suất (dạng cổ điển):}
    \[
    A = \{NN,\ NS\} \implies P(A) = \frac{2}{4} = \frac{1}{2};
    \]
    \[
    B = \{NN,\ SN\} \implies P(B) = \frac{2}{4} = \frac{1}{2};
    \]
    \[
    C = \{NN,\ SS\} \implies P(C) = \frac{2}{4} = \frac{1}{2}.
    \]
    \item \textbf{Chứng minh $A, B, C$ đôi một độc lập.} Theo bài giảng, $A$ và $B$ độc lập khi $P(A|B) = P(A)$, tương đương với $P(AB) = P(A)\,P(B)$ (do công thức nhân $P(AB) = P(B)\,P(A|B)$). Ta kiểm tra từng cặp:
    \begin{itemize}
        \item $AB = \{NN\} \implies P(AB) = \dfrac{1}{4} = \dfrac{1}{2} \cdot \dfrac{1}{2} = P(A)\,P(B)$. Vậy $A, B$ độc lập.
        \item $AC = \{NN\} \implies P(AC) = \dfrac{1}{4} = \dfrac{1}{2} \cdot \dfrac{1}{2} = P(A)\,P(C)$. Vậy $A, C$ độc lập.
        \item $BC = \{NN\} \implies P(BC) = \dfrac{1}{4} = \dfrac{1}{2} \cdot \dfrac{1}{2} = P(B)\,P(C)$. Vậy $B, C$ độc lập.
    \end{itemize}
    Chẳng hạn với cặp $A, C$: $P(A|C) = \dfrac{P(AC)}{P(C)} = \dfrac{1/4}{1/2} = \dfrac{1}{2} = P(A)$, nghĩa là việc biết $C$ xảy ra không làm thay đổi xác suất của $A$.

    Vậy $A, B, C$ đôi một độc lập. \hfill $\square$
\end{itemize}
\noindent\textbf{Nhận xét.} $ABC = \{NN\}$ nên $P(ABC) = \dfrac{1}{4} \ne \dfrac{1}{8} = P(A)\,P(B)\,P(C)$. Do đó $A, B, C$ đôi một độc lập nhưng \textit{không} độc lập toàn phần: nếu đã biết $A$ và $B$ cùng xảy ra thì chắc chắn $C$ xảy ra.

\vspace{1em}
%==================================================================
\begin{cau}
Có 2 lô hàng: lô 1 có 10 sản phẩm loại A, 5 sản phẩm loại B; lô 2 có 12 sản phẩm loại A, 8 sản phẩm loại B. Chọn ngẫu nhiên từ mỗi lô một sản phẩm. Tính xác suất chọn được:
\begin{enumerate}[label=\arabic*.]
    \item Hai sản phẩm loại A.
    \item Hai sản phẩm cùng loại.
    \item Hai sản phẩm khác loại.
\end{enumerate}
\end{cau}

\giai
Gọi $A_i$: ``Sản phẩm chọn từ lô $i$ là loại A''; $B_i = \overline{A}_i$: ``Sản phẩm chọn từ lô $i$ là loại B'' ($i = 1, 2$).
\[
P(A_1) = \frac{10}{15} = \frac{2}{3},\quad P(B_1) = \frac{5}{15} = \frac{1}{3}; \qquad
P(A_2) = \frac{12}{20} = \frac{3}{5},\quad P(B_2) = \frac{8}{20} = \frac{2}{5}.
\]
Việc chọn ở hai lô không ảnh hưởng lẫn nhau nên $A_1$ và $A_2$ độc lập; theo bài giảng, các cặp $(A_1, B_2)$, $(B_1, A_2)$, $(B_1, B_2)$ cũng độc lập.
\begin{enumerate}[label=\arabic*.]
    \item Hai sản phẩm loại A là biến cố $A_1 A_2$. Theo công thức nhân cho hai biến cố độc lập:
    \[
    P(A_1 A_2) = P(A_1)\,P(A_2) = \frac{2}{3} \cdot \frac{3}{5} = \frac{2}{5} = 0{,}4.
    \]
    \item Gọi $C$: ``Hai sản phẩm cùng loại'' $\implies C = A_1 A_2 \cup B_1 B_2$, với $A_1 A_2$ và $B_1 B_2$ xung khắc. Theo công thức cộng và công thức nhân:
    \[
    P(C) = P(A_1)\,P(A_2) + P(B_1)\,P(B_2) = \frac{2}{3} \cdot \frac{3}{5} + \frac{1}{3} \cdot \frac{2}{5} = \frac{6}{15} + \frac{2}{15} = \frac{8}{15} \approx 0{,}5333.
    \]
    \item ``Hai sản phẩm khác loại'' là biến cố đối lập $\overline{C}$:
    \[
    P(\overline{C}) = 1 - P(C) = 1 - \frac{8}{15} = \frac{7}{15} \approx 0{,}4667.
    \]
    \textit{Kiểm tra:} $\overline{C} = A_1 B_2 \cup B_1 A_2 \implies P(\overline{C}) = \dfrac{2}{3} \cdot \dfrac{2}{5} + \dfrac{1}{3} \cdot \dfrac{3}{5} = \dfrac{4}{15} + \dfrac{3}{15} = \dfrac{7}{15}$.
\end{enumerate}
\dapso{1. $\dfrac{2}{5} = 0{,}4$;\quad 2. $\dfrac{8}{15} \approx 0{,}5333$;\quad 3. $\dfrac{7}{15} \approx 0{,}4667$.}

\vspace{1em}
%==================================================================
\begin{cau}
Một lô hàng gồm 20 sản phẩm trong đó có 2 phế phẩm. Người ta lần lượt lấy mỗi lần một sản phẩm để kiểm tra (không hoàn lại) cho đến khi phát hiện đủ 2 phế phẩm thì dừng lại.
\begin{enumerate}[label=\alph*/]
    \item Tính xác suất để việc kiểm tra dừng lại ở lần thứ hai.
    \item Tính xác suất để việc kiểm tra dừng lại ở lần thứ ba.
\end{enumerate}
\end{cau}

\giai
Lô hàng có 2 phế phẩm và 18 chính phẩm. Gọi $F_i$: ``Sản phẩm kiểm tra lần thứ $i$ là phế phẩm'' ($i = 1, 2, \dots$). Vì lấy không hoàn lại nên các biến cố này không độc lập; ta dùng công thức nhân cho $n$ biến cố tùy ý (tương tự Ví dụ về thử bóng đèn trong bài giảng).
\begin{enumerate}[label=\alph*/]
    \item Gọi $H_2$: ``Việc kiểm tra dừng lại ở lần thứ hai''. Điều này xảy ra khi và chỉ khi cả hai lần đầu đều lấy phải phế phẩm: $H_2 = F_1 F_2$. Do đó
    \[
    P(H_2) = P(F_1)\,P(F_2|F_1) = \frac{2}{20} \cdot \frac{1}{19} = \frac{1}{190} \approx 0{,}0053.
    \]
    \item Gọi $H_3$: ``Việc kiểm tra dừng lại ở lần thứ ba''. Điều này xảy ra khi trong 2 lần đầu có đúng 1 phế phẩm và lần thứ ba lấy được phế phẩm còn lại:
    \[
    H_3 = F_1 \overline{F}_2 F_3 \ \cup\ \overline{F}_1 F_2 F_3,
    \]
    hai biến cố này xung khắc. Áp dụng công thức cộng và công thức nhân:
    \begin{align*}
    P(H_3) &= P(F_1)\,P(\overline{F}_2|F_1)\,P(F_3|F_1\overline{F}_2) + P(\overline{F}_1)\,P(F_2|\overline{F}_1)\,P(F_3|\overline{F}_1 F_2) \\
           &= \frac{2}{20} \cdot \frac{18}{19} \cdot \frac{1}{18} + \frac{18}{20} \cdot \frac{2}{19} \cdot \frac{1}{18} \\
           &= \frac{1}{190} + \frac{1}{190} = \frac{1}{95} \approx 0{,}0105.
    \end{align*}
\end{enumerate}
\dapso{a/ $\dfrac{1}{190} \approx 0{,}0053$;\quad b/ $\dfrac{1}{95} \approx 0{,}0105$.}

\vspace{1em}
%==================================================================
\begin{cau}
Bắn 3 viên đạn độc lập vào một bia. Xác suất trúng bia của mỗi viên đạn lần lượt là $0{,}6$; $0{,}9$; $0{,}7$. Tính xác suất:
\begin{enumerate}[label=\alph*/]
    \item Cả ba viên đều trúng bia.
    \item Không viên nào trúng bia.
    \item Có đúng một viên trúng bia.
    \item Có ít nhất 1 viên trúng bia.
\end{enumerate}
\end{cau}

\giai
Gọi $A_i$: ``Viên đạn thứ $i$ trúng bia'' ($i = 1, 2, 3$). Theo giả thiết:
\[
P(A_1) = 0{,}6,\ P(\overline{A}_1) = 0{,}4;\quad
P(A_2) = 0{,}9,\ P(\overline{A}_2) = 0{,}1;\quad
P(A_3) = 0{,}7,\ P(\overline{A}_3) = 0{,}3.
\]
Các viên đạn được bắn độc lập nên $A_1, A_2, A_3$ độc lập toàn phần (và các biến cố đối lập của chúng cũng vậy). Ta dùng công thức nhân cho các biến cố độc lập.
\begin{enumerate}[label=\alph*/]
    \item Cả ba viên trúng: $A_1 A_2 A_3$.
    \[
    P(A_1 A_2 A_3) = P(A_1)\,P(A_2)\,P(A_3) = 0{,}6 \cdot 0{,}9 \cdot 0{,}7 = 0{,}378.
    \]
    \item Không viên nào trúng: $\overline{A}_1 \overline{A}_2 \overline{A}_3$.
    \[
    P(\overline{A}_1 \overline{A}_2 \overline{A}_3) = 0{,}4 \cdot 0{,}1 \cdot 0{,}3 = 0{,}012.
    \]
    \item Có đúng một viên trúng:
    \[
    H_1 = A_1 \overline{A}_2 \overline{A}_3 \cup \overline{A}_1 A_2 \overline{A}_3 \cup \overline{A}_1 \overline{A}_2 A_3,
    \]
    ba biến cố này xung khắc từng đôi, nên
    \begin{align*}
    P(H_1) &= 0{,}6 \cdot 0{,}1 \cdot 0{,}3 + 0{,}4 \cdot 0{,}9 \cdot 0{,}3 + 0{,}4 \cdot 0{,}1 \cdot 0{,}7 \\
           &= 0{,}018 + 0{,}108 + 0{,}028 = 0{,}154.
    \end{align*}
    \item Có ít nhất 1 viên trúng: $A_1 \cup A_2 \cup A_3$, là biến cố đối lập của ``không viên nào trúng'':
    \[
    P(A_1 \cup A_2 \cup A_3) = 1 - P(\overline{A}_1 \overline{A}_2 \overline{A}_3) = 1 - 0{,}012 = 0{,}988.
    \]
\end{enumerate}
\dapso{a/ $0{,}378$;\quad b/ $0{,}012$;\quad c/ $0{,}154$;\quad d/ $0{,}988$.}

\vspace{1em}
%==================================================================
\begin{cau}
Hộp thứ nhất chứa 3 bi đỏ và 2 bi xanh; hộp thứ hai chứa 2 bi đỏ và 8 bi xanh. Chọn ngẫu nhiên một hộp và từ hộp đó chọn ngẫu nhiên một viên bi.
\begin{enumerate}[label=\arabic*.]
    \item Tìm xác suất chọn được bi đỏ.
    \item Giả sử đã chọn được bi đỏ. Tính xác suất để bi đỏ này là của hộp thứ nhất.
\end{enumerate}
\end{cau}

\giai
Gọi $H_i$: ``Chọn được hộp thứ $i$'' ($i = 1, 2$). Hộp được chọn ngẫu nhiên nên
\[
P(H_1) = P(H_2) = \frac{1}{2},
\]
và $\{H_1, H_2\}$ là một hệ đầy đủ. Gọi $D$: ``Viên bi chọn được màu đỏ''. Ta có
\[
P(D|H_1) = \frac{3}{5} = 0{,}6; \qquad P(D|H_2) = \frac{2}{10} = 0{,}2.
\]
\begin{enumerate}[label=\arabic*.]
    \item Bài toán hỏi xác suất của biến cố chung $D$, nên dùng công thức xác suất đầy đủ:
    \[
    P(D) = P(H_1)\,P(D|H_1) + P(H_2)\,P(D|H_2) = \frac{1}{2} \cdot 0{,}6 + \frac{1}{2} \cdot 0{,}2 = 0{,}3 + 0{,}1 = 0{,}4.
    \]
    \item Biết kết quả $D$ đã xảy ra, cần tìm xác suất của nguyên nhân $H_1$, nên dùng công thức Bayes:
    \[
    P(H_1|D) = \frac{P(H_1)\,P(D|H_1)}{P(D)} = \frac{0{,}3}{0{,}4} = \frac{3}{4} = 0{,}75.
    \]
\end{enumerate}
\noindent\textit{Sơ đồ cây:} Nhánh hộp 1 -- đỏ: $0{,}5 \cdot 0{,}6 = 0{,}3$; nhánh hộp 2 -- đỏ: $0{,}5 \cdot 0{,}2 = 0{,}1$.
$P(D) = 0{,}3 + 0{,}1 = 0{,}4$; \ $P(H_1|D) = \dfrac{0{,}3}{0{,}4} = 0{,}75$.

\dapso{1. $0{,}4$;\quad 2. $0{,}75$.}

\vspace{1em}
%==================================================================
\begin{cau}
Một bác sĩ chữa khỏi bệnh D cho một người với xác suất là 95\%. Giả sử có 10 người bị bệnh D đến bác sĩ này chữa bệnh, tính xác suất:
\begin{enumerate}[label=\alph*)]
    \item Có 8 người khỏi bệnh.
    \item Có nhiều nhất 9 người khỏi bệnh.
\end{enumerate}
\end{cau}

\giai
Việc chữa bệnh cho 10 người là 10 phép thử độc lập. Trong mỗi phép thử chỉ có hai kết quả: khỏi bệnh (biến cố $A$) hoặc không khỏi ($\overline{A}$), với $P(A) = p = 0{,}95$ không đổi. Đây là dãy phép thử Bernoulli với
\[
n = 10, \qquad p = 0{,}95, \qquad q = 1 - p = 0{,}05.
\]
\begin{enumerate}[label=\alph*)]
    \item Xác suất có đúng 8 người khỏi bệnh ($k = 8$):
    \[
    P_{10}(8;\ 0{,}95) = C_{10}^8\,(0{,}95)^8\,(0{,}05)^2 = 45 \cdot 0{,}66342 \cdot 0{,}0025 \approx 0{,}0746.
    \]
    \item ``Có nhiều nhất 9 người khỏi bệnh'' nghĩa là $k \in \{0, 1, \dots, 9\}$. Biến cố đối lập là ``Cả 10 người đều khỏi bệnh'' ($k = 10$):
    \[
    P_{10}(10;\ 0{,}95) = C_{10}^{10}\,(0{,}95)^{10}\,(0{,}05)^0 = (0{,}95)^{10} \approx 0{,}5987.
    \]
    Vậy
    \[
    P = 1 - P_{10}(10;\ 0{,}95) = 1 - (0{,}95)^{10} \approx 1 - 0{,}5987 = 0{,}4013.
    \]
\end{enumerate}
\dapso{a) $\approx 0{,}0746 = 7{,}46\%$;\quad b) $\approx 0{,}4013 = 40{,}13\%$.}

\vspace{1em}
%==================================================================
\begin{cau}
Xét việc gieo một con súc sắc cân đối và đồng chất. Hãy tính xem cần gieo bao nhiêu lần thì xác suất được ít nhất 1 lần xuất hiện mặt sáu chấm không nhỏ hơn $0{,}9$.
\end{cau}

\giai
Gieo súc sắc $n$ lần là dãy $n$ phép thử Bernoulli. Trong mỗi lần gieo, biến cố $A$: ``Xuất hiện mặt sáu chấm'' có xác suất
\[
p = \frac{1}{6}, \qquad q = 1 - p = \frac{5}{6}.
\]
Biến cố ``Có ít nhất 1 lần xuất hiện mặt sáu chấm'' là đối lập của biến cố ``Không lần nào xuất hiện mặt sáu chấm'' ($k = 0$), nên
\[
P = 1 - P_n(0;\ \tfrac{1}{6}) = 1 - C_n^0 \left(\frac{1}{6}\right)^0 \left(\frac{5}{6}\right)^n = 1 - \left(\frac{5}{6}\right)^n.
\]
Theo yêu cầu đề bài:
\begin{align*}
1 - \left(\frac{5}{6}\right)^n \ge 0{,}9
&\iff \left(\frac{5}{6}\right)^n \le 0{,}1 \\
&\iff n \ge \log_{5/6}(0{,}1) = \frac{\ln 0{,}1}{\ln (5/6)} \approx \frac{-2{,}3026}{-0{,}1823} \approx 12{,}63
\end{align*}
(bất đẳng thức đổi chiều vì cơ số $\dfrac{5}{6} < 1$).

Vì $n$ là số nguyên dương nên $n_{\min} = 13$.

\textit{Kiểm tra:} $n = 12$: $1 - (5/6)^{12} \approx 0{,}8878 < 0{,}9$; \ $n = 13$: $1 - (5/6)^{13} \approx 0{,}9065 \ge 0{,}9$.

\dapso{Cần gieo ít nhất $13$ lần.}

\end{document}
